\section{An explicit bound for every fixed Eisenstein Weil fourfold family}
\label{sec:sixfold-extension}

The principal fourfold bound in R10--R12 extends to arbitrary fixed determinant and polarization data by a quantitative version of Schoen's compensation argument \cite[Proposition 10]{Schoen1998}. The intermediate variety has dimension six. We give the constants, the rational lattice algorithm, and the cost of each isogeny explicitly. The resulting integer is an upper bound; it is not the least complexity at an unspecified member of the family.

Throughout, $K=\Q(\omega)$ with $\omega^2+\omega+1=0$, and $\mathcal O_K=\Z[\omega]$. Polarizations in the main statement are compatible with the $K$ action, meaning that their Rosati involution induces complex conjugation on $K$. A nonzero Weil class on a fourfold has positive square $T=t^2$.

\subsection{R20: a uniformly bounded detecting cycle on principal sixfolds}

Define
\begin{equation}
D_6=\frac{110224800\cdot14171760}{27}
=57\,854\,793\,024\,000.
\label{eq:D6}
\end{equation}

\begin{proposition}[R20]
\label{prop:R20}
In a Zariski dense subset of the principally polarized Eisenstein Weil sixfold component of signature $(3,3)$, there is an effective three-dimensional cycle $Z$ of polarization degree at most $D_6$ whose projection to the Weil plane is nonzero.
\end{proposition}

\begin{proof}
Take a connected unramified cyclic triple cover $\pi:C\to X$ with $g(X)=4$, so $g(C)=10$. Write $J=\operatorname{Jac}(C)$, let $\Theta$ be its principal polarization, and let $B\subset J$ be the connected Prym, of dimension six. Set $L=\Theta|_B$.

We first determine its polarization, including the connected-kernel issue. Let $Q_0=\im(\pi^*:\operatorname{Jac}(X)\to J)$, and let $\eta\in\operatorname{Pic}^0(X)[3]$ define the connected cover. The eigensheaf decomposition
\[
\pi_*\mathcal O_C=\mathcal O_X\oplus\eta^{-1}\oplus\eta^{-2}
\]
shows that $\pi^*M$ is trivial for $M\in\operatorname{Pic}^0(X)$ exactly when $M\in\langle\eta\rangle$: if the pullback is trivial it has a nonzero section, and the projection formula forces one of the degree-zero bundles $M\otimes\eta^{-i}$ to be trivial; the converse follows from the cover construction. Thus $\operatorname{Jac}(X)\to Q_0$ has degree three. The dual of pullback, under the canonical principal polarizations, is the norm, and $\operatorname{Nm}\circ\pi^*=[3]$. Consequently $(\pi^*)^*(\Theta|_{Q_0})=3\Theta_X$ as polarization classes. By \cite[Theorem 13.3, p.~25; Remark 16.9, p.~33]{MilneAV}, the polarization on $Q_0$ has degree
\[
\deg\lambda_{Q_0}=3^8/3^2=3^6.
\]
Here $B=\ker(\operatorname{Nm})^0=Q_0^\perp$. Indeed, the restriction map $J^\vee\to Q_0^\vee$ has connected kernel, whereas the dual of the degree-three map $\operatorname{Jac}(X)\to Q_0$ is an isogeny; their composite is the norm under the principal-polarization identifications. Dualizing $0\to Q_0\to J\to J/Q_0\to0$ and its complementary sequence gives $B^\perp=Q_0$. The kernels of the two restricted polarizations are therefore the same scheme-theoretic intersection $B\cap Q_0$; this is the dual-exact-sequence argument of \cite[\S11, pp.~17--18]{MilneAV}. In particular $K(L)$ has order $3^6$, with no assumption that the entire norm kernel is connected. The operator $1+\sigma+\sigma^2$ is zero on $B$ and multiplication by three on $Q_0$, so $K(L)$ is killed by three. Its type is therefore
\[
(1,1,1,3,3,3),\qquad L^6=6!\cdot27=19440.
\]
The deck transformation is the identity on $Q_0$, hence on $K(L)$. Every maximal isotropic subgroup $H\subset K(L)$, of order $27$, is consequently $\mathcal O_K$-stable. The quotient
\[
q:B\longrightarrow Q=B/H
\]
has a polarization $\theta$ with $q^*\theta=L$ by \cite[Proposition 16.8, p.~33]{MilneAV}. The degree identity in Remark~16.9 there gives $\deg\lambda_\theta=3^6/27^2=1$, so it is principal. The $\mathcal O_K$ action descends, and remains compatible with $\theta$ because it preserves $L$ and isogeny pullback is injective on rational divisor classes.

For the detecting cycle, use the $r=0$ case of \cite[Theorem 2.0 and its proof, pp.~12--13]{Schoen1988}. Its cycles are supported over
\[
|K_X|=\Pp^3\subset\operatorname{Sym}^6X.
\]
The maps $C^6\to\operatorname{Sym}^6C\to\operatorname{Sym}^6X$ are finite. In the notation of that construction, the two-dimensional rational space $U$ is the pullback of the Weil plane of $B$ under
\[
f=P\circ\operatorname{Abel}:C^6\longrightarrow B,
\qquad P=3\id-(1+\sigma+\sigma^2).
\]
This map is dominant and generically finite. The Abel image has class $[W_6(C)]=\Theta^4/4!$ by Poincar\'e's formula \cite[p.~149, formulas (6)--(7)]{GriffithsPoincare1982}. Its intersection number with a translate of $Q_0=\ker(P)^0$ is $\int_{Q_0}\Theta^4/4!>0$. Thus a general fiber of $P$ meets $W_6(C)$, proving dominance; the source and target both have dimension six. Consequently $f_*f^*$ is multiplication by a nonzero degree. The identification of $U$ with this pullback is \cite[Remark 1.4, pp.~6--7; \S3, p.~24]{Schoen1988}. Some constituent of the detecting cycles supported above $|K_X|$ has a pushforward with nonzero Weil projection.

Its Abel image is contained in $W_6(C)\cap T$, for a translate $T$ of $B$ in the norm fiber of $K_X$. Every point of this intersection is represented by an effective divisor whose norm is canonical. The finite symmetric-product map shows that the intersection has dimension at most three. A contributing three-dimensional image is therefore an irreducible component $S$ of $W_6(C)\cap T$. On $T$, the map $P$ is multiplication by three followed by translation. This acts invertibly on rational cohomology, so the reduced component $S$, translated into $B$, itself has nonzero Weil projection. Selecting the reduced image eliminates any need to charge an ordered-product multiplicity or the degree of $[3]$.

Embed $J$ by $3\Theta$, which is very ample by \cite[Theorem 7.3, p.~11]{MilneAV}. Poincar\'e's formula gives
\[
\deg_{3\Theta}W_6(C)=3^6\frac{10!}{4!}=110224800,
\qquad
\deg_{3\Theta}B=3^6\cdot19440=14171760.
\]
The general projective B\'ezout inequality bounds the degree of every reduced irreducible component of the intersection by the product of these two degrees, even when the intersection is improper \cite[Example 8.4.6]{Fulton}. Since $S$ has dimension three,
\[
\deg_LS\leq\frac{110224800\cdot14171760}{3^3}=D_6.
\]
Finally $q_*S$ is effective, has nonzero Weil projection, and has $\theta$-degree equal to $\deg_LS$ by the projection formula.

It remains to justify density in the claimed component. Sections 11--12 of \cite[pp.~333--334]{Schoen1998} construct precisely this genus-four cyclic-cover Prym family, and Lemma 13 on pp.~334--335 proves that its period map to the full corresponding Weil-pair component is dominant. For clarity, the decisive finite calculation in that proof can be checked directly. At its exponent-three cover with group $(\Z/3)^3$, the two nontrivial quotient-character spaces of holomorphic differentials have character triples
\[
\{(2,2,0),(2,2,1),(2,2,2)\},\qquad
\{(0,2,2),(1,1,2),(2,0,2)\}.
\]
All nine pairwise sums are distinct. The factors are one-dimensional, and the product of nonzero differentials on an integral curve is nonzero. Thus the nine multiplication maps land injectively in distinct eigenspaces, proving the injectivity of the cotangent map used by Schoen.

The isotropic quotient is a finite isogeny correspondence, so it preserves dimension and gives a dense image in a principal component. A unimodular Eisenstein Hermitian lattice of signature $(3,3)$ is $U^{\oplus3}$: apply Meyer's theorem to obtain a primitive isotropic vector, use unimodularity to find a partner of pairing one, and use $\Tr(\mathcal O_K)=\Z$ to make that partner isotropic. This splits a hyperbolic plane integrally. Repeat in signature $(2,2)$. The final binary lattice has determinant $-1$, is rationally isotropic, and splits by the same primitive-vector argument. This is the dimension-six version of the component check in R12; Meyer's theorem is \cite[Corollary 3.3.4]{PetersSterk}. Hence the principal component is the one asserted in the proposition.
\end{proof}

\subsection{R21: explicit extraction on a principal sixfold}

For $\delta\geq1$ and $U>0$, put
\begin{equation}
C_6(\delta,U)=516\,560\,652\,\delta D_6^4
\left(100+1029\left\lceil\sqrt U\right\rceil\right).
\label{eq:C6}
\end{equation}

\begin{theorem}[R21]
\label{thm:R21}
Let $(Q,\theta)$ be a principally polarized Eisenstein Weil sixfold of signature $(3,3)$, and let $\tau\ne0$ belong to its rational Weil plane. Suppose $\delta\tau$ is integral, and write
\[
U=-\int_Q\tau^2>0.
\]
There are effective integral three-cycles $Z^+,Z^-$ and an integer $N\geq1$ with
\[
[Z^+]-[Z^-]=N\tau,\qquad
N,\deg_\theta Z^+,\deg_\theta Z^-\leq C_6(\delta,U).
\]
\end{theorem}

\begin{proof}
First work on the dense subset of Proposition~\ref{prop:R20}. Write
\[
R=[\omega]^*,\qquad S=[\sqrt{-3}]^*
\]
on $H^6(Q,\Q)$. On the summand with $a$ factors in the first $K$-eigenspace and $6-a$ in the second, these have eigenvalues $\omega^{-a}$ and $-27(-1)^a$. Therefore
\begin{equation}
162e_W=(27-S)(1+R+R^2)
\label{eq:sixfold-projector}
\end{equation}
is $162$ times the projector to the two Weil lines $a=0,6$.

Choose a detecting effective cycle $Z$ of degree at most $D_6$ and put $u=162e_W[Z]\ne0$. The formula
\[
u=27\bigl([Z]+R[Z]+R^2[Z]\bigr)
-S\bigl([Z]+R[Z]+R^2[Z]\bigr)
\]
is an integral signed-cycle representative. The unit maps preserve $\theta$. Pullback by $[\sqrt{-3}]$ multiplies codimension-three degree by $27$, since its pullback polarization is $3\theta$. Thus both numerator degrees are at most $81D_6$.

The Weil plane is primitive of Hodge type $(3,3)$, so its positive form is
\[
\langle x,y\rangle=-\int_Qxy.
\]
The sign is essential in codimension three. The same flat-metric argument as R11 has total volume $\theta^6/6!=1$, and a complex three-cycle of degree $D$ has volume $D/6$. A harmonic form is translation-invariant, and its comass is at most its pointwise Euclidean norm. Consequently
\[
\|[Z]\|\leq D_6/6.
\]
The $K$-weight projectors are orthogonal because the Rosati involution is conjugation; hence, with $d_u=\langle u,u\rangle$,
\begin{equation}
0<d_u\leq162^2(D_6/6)^2=729D_6^2=:M_6.
\label{eq:M6}
\end{equation}
The integer $d_u$ is the positive primitive intersection norm.

Here $R$ acts trivially on the Weil plane, so the fourfold basis $u,Ru$ cannot be reused. Instead take
\[
\alpha=2-\omega,\qquad T_\alpha=[\alpha]^*.
\]
One computes
\[
\alpha^3=1-18\omega,\quad
\alpha^6=-323-360\omega,\quad
\Re(\alpha^6)=-143,\quad |\alpha^6|=343.
\]
The value $\alpha^6$ is nonreal, so $u,T_\alpha u$ form a rational basis of $W$. Their positive Gram matrix is
\[
d_u\begin{pmatrix}1&-143\\-143&117649\end{pmatrix},
\qquad \Delta=97200d_u^2.
\]
The signed representative of $T_\alpha u$ has numerator degrees at most $343\cdot81D_6$.

The integers
\[
p=\langle\delta\tau,u\rangle,\qquad
q=\langle\delta\tau,T_\alpha u\rangle
\]
satisfy, by Cramer's rule,
\[
\Delta\delta\tau=a u+bT_\alpha u,
\quad
a=d_u(117649p+143q),\qquad
b=d_u(143p+q).
\]
Cauchy--Schwarz in this positive two-dimensional space yields
\[
|p|\leq\delta\sqrt{Ud_u},\qquad
|q|\leq343\delta\sqrt{Ud_u},
\]
and therefore
\[
|a|\leq166698\delta\sqrt U\,d_u^{3/2},\qquad
|b|\leq486\delta\sqrt U\,d_u^{3/2}.
\]
Combining the signed representatives gives
\begin{align*}
N&=\delta\Delta\leq97200\delta M_6^2,\\
\deg Z^\pm&\leq81D_6(|a|+343|b|)
\leq27005076\delta D_6\sqrt U\,M_6^{3/2}.
\end{align*}
Since $M_6^{3/2}=19683D_6^3$, the two coefficients are respectively
\[
97200\cdot729^2=51\,656\,065\,200,
\quad
27005076\cdot19683=531\,540\,910\,908.
\]
Their sum, with $\lceil\sqrt U\rceil$ in the second term, is exactly \eqref{eq:C6} and bounds both quantities.

Finally pass to a finite level cover trivializing the Weil monodromy, as in R13a. Its determinant character takes values in $\mathcal O_K^\times=\mu_6$, so a cover of degree at most six suffices. Fix $\tau$ on that cover. The integers $\delta$ and the norm $U$ are constant. The bounded-class locus of R13b is closed, contains the dense Prym-quotient image, and therefore is the whole component. This proves the bound at every point, not just a general point.
\end{proof}

\subsection{R22: a constructive rational isometry and its exact cost}

\begin{proposition}[R22]
\label{prop:R22}
Let $(A,H)$ be an Eisenstein Weil fourfold of signature $(2,2)$ with integral $\mathcal O_K$ action. Set
\[
d=\chi(H)=H^4/4!\in\Z_{>0}.
\]
There is an explicit terminating rational-linear-algebra procedure, using a search for isotropic vectors, which produces a rational Hermitian isometry and an integer $m\geq1$ defining an isogeny
\[
h:A\times E^2\longrightarrow Q
\]
to a principally polarized Eisenstein Weil sixfold. On $E^2$ use the mixed $K$ action and $H'=\theta_1+d\theta_2$, and put $H_0=H+H'$. Then
\begin{equation}
h^*\theta_Q=m^2H_0,\qquad
e:=\deg h=m^{12}d^2.
\label{eq:comp-isogeny}
\end{equation}
The matrix, $m$, and $e$ depend only on the fixed polarized Weil lattice, not on the point of its period domain.
\end{proposition}

\begin{proof}
Let $\Lambda=H_1(A,\Z)$ and let $E_H$ be its alternating polarization form. Since $\mathcal O_K$ is a principal ideal domain, $\Lambda$ is free of rank four over $\mathcal O_K$. Define
\[
\mathfrak h_H(x,y)=E_H(\omega x,y)-\bar\omega E_H(x,y).
\]
This is an integral Hermitian form and
\[
E_H(x,y)=\Tr_{K/\Q}\left(\frac{\mathfrak h_H(x,y)}{\sqrt{-3}}\right).
\]
The inverse different in this trace formula makes the identity Hermitian lattice yield a unimodular alternating lattice. The determinant formula for trace pairings then gives
\[
\det_{\Z}(E_H)=(\det_K\mathfrak h_H)^2.
\]
The polarization Euler characteristic is the positive square root of the alternating determinant. Signature $(2,2)$ gives a positive Hermitian determinant, so
\[
\det_K\mathfrak h_H=\chi(H)=d.
\]

The auxiliary polarized surface $(E^2,H')$ has signature $(1,1)$, Hermitian determinant $-d$, and $\chi(H')=d$. Thus the product has signature $(3,3)$ and determinant $-d^2$, the determinant class of a hyperbolic six-dimensional Hermitian space. This is also the compensation in \cite[\S7 and Proposition 10]{Schoen1998}.

Here is a constructive version of the rational isometry, without assuming an uncomputed classification table. The underlying rational quadratic form has dimension twelve and signature $(6,6)$. Meyer's theorem \cite[Corollary 3.3.4]{PetersSterk} gives a nonzero rational isotropic vector. Enumerate rational vectors and test the exact norm equation until one is found. Nondegeneracy gives a rational partner of Hermitian pairing one; subtract half its rational norm times the isotropic vector to make it isotropic too. This splits a rational hyperbolic plane. In the remaining signature $(2,2)$ Hermitian space, apply the same procedure to its dimension-eight rational quadratic form and split a second plane.

Let $A_0$ denote the accumulated $K$-basis-change matrix. The remaining binary Hermitian form has determinant
\[
-d^2\Nm_{K/\Q}(\det A_0)=-\Nm_{K/\Q}(s),
\qquad s=d\det A_0.
\]
If a diagonal entry is zero, its corresponding vector is already isotropic. Otherwise a determinant-one shear gives the diagonal form
\[
\left\langle a,-\frac{\Nm(s)}a\right\rangle,
\qquad a\in\Q^\times.
\]
The vector $(s,a)$ is isotropic. Split the last hyperbolic plane as above. All steps after the two terminating isotropic-vector searches are exact rational linear algebra. This produces an isometry $\iota$ to the standard unimodular Hermitian lattice $U^{\oplus3}$ over $K$.

Choose $m>0$ clearing the entries of $\iota$ relative to the two integral lattices. For each admissible complex structure $J$ on the source, use $\iota J\iota^{-1}$ on the target. The map $m\iota$ is an integral holomorphic lattice map, hence an isogeny to a principally polarized abelian variety $Q$; see \cite[\S6]{Schoen1998} for this passage between universal families. Because $\iota$ preserves the Hermitian and alternating forms, $h^*\theta_Q=m^2H_0$. Therefore
\[
\deg h=\frac{(h^*\theta_Q)^6}{\theta_Q^6}
=m^{12}\chi(H_0)
=m^{12}\chi(H)\chi(H')=m^{12}d^2.
\]
The construction uses only the fixed lattice and its forms. No bound asserting that $m$ is small in terms of $d$ alone is needed or claimed.
\end{proof}

For later use, let $\Gamma$ be a rational degree-six class on the product and define
\[
\tau=e^{-1}h_*\Gamma.
\]
Isogeny pullback and pushforward satisfy $h^*h_*=e$ on rational cohomology, so $h^*\tau=\Gamma$. If $3\delta\Gamma$ is integral, then
\begin{equation}
\delta_\tau=3\delta e\quad\hbox{clears }\tau,
\qquad
-\tau^2=(-\Gamma^2)/e.
\label{eq:comp-class-data}
\end{equation}
Proper pushforward preserves integral cohomology, which proves the denominator assertion. For an effective three-cycle $Y$ on $Q$, finite flatness of the isogeny and the projection formula give
\begin{equation}
\deg_{H_0}(h^*Y)
=\frac{e}{m^6}\deg_{\theta_Q}Y,
\qquad \frac{e}{m^6}=m^6d^2\in\Z_{>0}.
\label{eq:comp-pullback-degree}
\end{equation}

\subsection{R23: contraction with no additional numerator degree cost}

\begin{proposition}[R23]
\label{prop:R23}
For every $t$ in the fourfold Weil plane there is a class $\Gamma$ in the sixfold Weil plane of $A\times E^2$ such that
\[
-\Gamma^2=t^2,
\]
and $3\delta\Gamma$ is integral whenever $\delta t$ is integral. Contraction against a fixed divisor on $E^2$ sends $\Gamma$ to $t$. If $N\Gamma$ has an effective-difference representative with both $H_0$-degrees at most $C$, then $Nt$ has such a representative with both $H$-degrees at most $C$.
\end{proposition}

\begin{proof}
On the mixed-action surface $E^2$ let $F_1,F_2$ be the coordinate fibers and let $G$ be the graph of the identity. Put
\[
D=[G]-[F_1]-[F_2]=D^+-D^-,\qquad
D^+=[G],\quad D^-=[F_1]+[F_2].
\]
This lies in its Weil plane. Its intersections are
\[
D^2=-2,\qquad D\cdot R'D=D\cdot(R')^2D=1,
\]
where $R'=[\omega]^*$ for the mixed action. Let $R_A$ be the corresponding action on the fourfold Weil plane; it satisfies $1+R_A+R_A^2=0$. Define
\[
\Pi=\frac{1+R_A\otimes R'+R_A^2\otimes(R')^2}{3},
\qquad \Gamma=-\Pi(t\otimes D).
\]
After splitting $K$, the average keeps exactly the two tensor products with the same $K$ embedding in the fourfold and surface factors. These are the two sixfold Weil lines. The formula also shows that $3\delta\Gamma$ is integral if $\delta t$ is integral.

Contraction against $D$ is
\begin{align*}
p_{A*}(\Gamma\cdot p_{E^2}^*D)
&=-\frac{-2t+R_At+R_A^2t}{3}=t.
\end{align*}
Since $R_A$ is an isometry of the positive fourfold Weil plane,
\[
t\cdot R_At=t\cdot R_A^2t=-t^2/2.
\]
The averaging projector is self-adjoint, so
\[
\Gamma^2
=(t\otimes D)\cdot\Pi(t\otimes D)
=\frac{-2t^2-t^2/2-t^2/2}{3}=-t^2.
\]

It remains to control effective intersections. For the product principal surface polarization, $2(\theta_1+\theta_2)-[G]$ is nef: its Hermitian matrix is $\left(\begin{smallmatrix}1&1\\1&1\end{smallmatrix}\right)$. Since $d\geq1$ and $H'=\theta_1+d\theta_2$, this gives the nef comparisons
\[
D^+\leq2H',\qquad D^-\leq H'.
\]
For an effective three-cycle $X$ on $A\times E^2$, translate $G,F_1,F_2$ generally on $E^2$ so that their pullbacks intersect every component of $X$ properly. The translation incidence argument makes these intersections effective, including their multiplicities; their classes are unchanged. Pushforward to $A$ is effective, possibly zero. Projection formula and nef comparison give
\begin{align*}
\deg_Hp_{A*}(X\cdot D^+)&\leq2\int_XH^2H'
\leq\frac23\deg_{H_0}X,\\
\deg_Hp_{A*}(X\cdot D^-)&\leq\int_XH^2H'
\leq\frac13\deg_{H_0}X.
\end{align*}
For the last inequalities, $3H^2H'$ is one nonnegative summand in $(H+H')^3$ on the effective cycle $X$.

Now expand the intersection of $X^+-X^-$ with $D^+-D^-$. Its positive side is the sum of the $X^+D^+$ and $X^-D^-$ terms, and its negative side consists of the other two terms. Each degree is at most $(2/3+1/3)C=C$. The represented class is $Nt$ by the contraction calculation, and no multiplier has been added.
\end{proof}

\subsection{R24: the explicit uniform fourfold bound}

\begin{theorem}[R24]
\label{thm:R24}
Fix a compatible polarized Eisenstein Weil lattice of signature $(2,2)$ with integral $\mathcal O_K$ action. Let
\[
d=\chi(H),\qquad e=m^{12}d^2,
\]
where $m$ is the explicitly constructible lattice denominator of Proposition~\ref{prop:R22}. For a nonzero rational Weil vector $t$, choose $\delta\geq1$ with $\delta t$ integral and put $T=t^2>0$. Then, at every member of this family,
\begin{equation}
c_H(t)\leq C_4(\delta,T;H,m)
:=\frac{e}{m^6}\,
C_6\left(3\delta e,\frac{T}{e}\right).
\label{eq:C4}
\end{equation}
The same upper bound holds for its compatible realizations at every geometric good reduction of every arithmetic member. For $t=0$, the trivial bound is $1$.
\end{theorem}

\begin{proof}
Construct $\Gamma$ by Proposition~\ref{prop:R23}. It has $3\delta\Gamma$ integral and negative square $T$. Transport it to $\tau=e^{-1}h_*\Gamma$ on the principal sixfold. Formula \eqref{eq:comp-class-data} gives the valid class data
\[
\delta_\tau=3\delta e,\qquad U_\tau=T/e.
\]
Theorem~\ref{thm:R21} produces a representative with multiplier and numerator degrees at most $C_6(3\delta e,T/e)$. Pull it back by the finite flat isogeny $h$. The multiplier is unchanged and \eqref{eq:comp-pullback-degree} multiplies numerator degrees by $e/m^6$. This integer is at least one, so it bounds the multiplier as well. Proposition~\ref{prop:R23} contracts to $t$ without further degree or multiplier cost. This proves \eqref{eq:C4}.

All lattice quantities are fixed in the family, and the norm and denominator of a flat Weil vector are fixed after the finite monodromy cover of R13a. For an arithmetic member, descend the resulting complex cycles by B05 and specialize their effective closures. The intersection degree and all cycle-class comparisons are preserved. This gives one rational cycle in all realizations at every geometric good reduction, including primes dividing the displayed multiplier. No prime-distribution input is required for this bound.
\end{proof}

\paragraph{A rational action and its preliminary isogeny.}
If the original lattice $\Lambda$ is stable only under an order in $K$, choose an integer $f>0$ clearing the rational action of $\omega$ and set
\[
\widetilde\Lambda=f(\Lambda+\omega\Lambda)\subset\Lambda.
\]
This is $\mathcal O_K$-stable and has finite index $e_0=[\Lambda:\widetilde\Lambda]$. It defines an isogeny
\[
q:\widetilde A\longrightarrow A,
\qquad \widetilde H=q^*H,\quad\widetilde t=q^*t.
\]
The denominator $\delta$ still clears $\widetilde t$, while
\[
\chi(\widetilde H)=e_0\chi(H),\qquad
\widetilde t^2=e_0T.
\]
Apply Theorem~\ref{thm:R24} to these integral-action data, using its own fixed lattice denominator $\widetilde m$. Pushing cycles forward preserves their polarization degrees and changes the multiplier $N$ to $e_0N$. Therefore the fully charged extension is
\begin{equation}
c_H(t)\leq e_0\,
C_4(\delta,e_0T;\widetilde H,\widetilde m).
\label{eq:rational-action-bound}
\end{equation}
The direction of this preliminary isogeny matters: the stable \emph{sublattice} gives a cover of $A$, whose effective cycles can be pushed forward.

If a given polarization is not initially compatible, after integralizing the action replace it by
\[
\overline H=H+[\omega]^*H+[\omega^2]^*H.
\]
It is compatible and $\overline H-H$ is nef. Effective surface degrees for $H$ are at most those for $\overline H$, so the corresponding bound also bounds the original degrees.

These constructions complete the uniform bounded-degree and bounded-multiplier gate for each fixed T5 fourfold family and each fixed rational Weil vector. Their inputs include the actual lattice, polarization and class normalization. They do not give a least-complexity formula at an arbitrary point, a constant independent of all such inputs, or a solution for unrelated higher-dimensional Weil families.
